% ============================================================================
%  01_proposito.tex — IEEE 730 §1: Propósito y Alcance (Sección A.a)
% ============================================================================

\chapter{Propósito y Alcance}
\label{ch:proposito}

El presente \textbf{Plan de Aseguramiento de la Calidad de Software} (SQAP) establece la organización, las tareas,
las responsabilidades y los criterios de calidad para el proyecto \textbf{SpeechNotes --- Sistema de Transcripción
Inteligente con IA}. El documento sigue los lineamientos del estándar \textbf{IEEE 730-1998} adaptados al marco
académico del proyecto final de Aseguramiento de la Calidad de Software (SQA).

Este SQAP es ejecutado por un equipo externo e independiente al equipo de desarrollo, durante un \textit{sprint SQA}
comprendido entre el \textbf{18 y el 21 de julio de 2026}.

% ---------------------------------------------------------------------------
\section{Identificación del Sistema}
\label{sec:identificacion}

\begin{itemize}
    \item \textbf{Nombre del sistema:} SpeechNotes
    \item \textbf{Versión evaluada:} 1.0 (rama \texttt{planning-SQA})
    \item \textbf{Repositorio:} \url{https://github.com/anomalyco/SpeechNotes}
    \item \textbf{Stack tecnológico:} Python/FastAPI (backend), Next.js 16 (frontend),
          MongoDB/SQLite/ChromaDB (datos), NVIDIA NIM (modelos ASR, BNR, LLM), Socket.IO (tiempo real).
    \item \textbf{Arquitectura:} Monorepo con backend FastAPI + Socket.IO, frontend Next.js 16,
          apps de escritorio Tauri/Electron.
\end{itemize}

% ---------------------------------------------------------------------------
\section{Alcance: Módulos Incluidos}
\label{sec:modulos-incluidos}

El presente plan cubre los siguientes \textbf{10 requisitos funcionales}, seleccionados por representar el flujo
crítico de valor del usuario (registro de audio → transcripción → gestión → inteligencia artificial):

\begin{longtable}{clp{7cm}}
\caption{Módulos incluidos en el alcance del SQAP}
\label{tab:modulos-incluidos}\\
\toprule
\textbf{RF} & \textbf{Módulo} & \textbf{Justificación de inclusión}\\
\midrule
\endfirsthead
\toprule
\textbf{RF} & \textbf{Módulo} & \textbf{Justificación de inclusión}\\
\midrule
\endhead
\bottomrule
\endfoot

RF-01 & Auth (Autenticación) & Punto de entrada crítico: sin autenticación ningún otro módulo es accesible. Un fallo aquí bloquea el 100\% de los casos de prueba.\\

RF-02 & Audio / Frontend (Micrófono) & Módulo de captura de audio en el navegador. Funcionalidad principal del producto.\\

RF-03 & Audio / Backend (VAD) & Detección de actividad de voz (Voice Activity Detection). Controla inicio y fin de grabación.\\

RF-04 & Realtime ASR & Transcripción en tiempo real vía Socket.IO y NVIDIA NIM (Parakeet/Canary).\\

RF-05 & Upload / Backend & Carga de archivos de audio existentes para transcribir asíncronamente.\\

RF-06 & Transcripciones (CRUD) & Listado, visualización y eliminación de notas de voz guardadas.\\

RF-07 & Editor & Edición y guardado de transcripciones por parte del usuario.\\

RF-08 & Búsqueda & Búsqueda textual sobre el contenido de las transcripciones guardadas.\\

RF-09 & IA / Formatter & Formateo y estructuración automática de texto crudo mediante un modelo LLM de NVIDIA NIM.\\

RF-10 & Chat / RAG & Chat contextual sobre las transcripciones mediante un agente RAG con ChromaDB y NVIDIA NIM.\\

\bottomrule
\end{longtable}

% ---------------------------------------------------------------------------
\section{Alcance: Módulos Excluidos}
\label{sec:modulos-excluidos}

Los siguientes módulos y actividades quedan \textbf{fuera del alcance} de este plan, con su justificación correspondiente:

\begin{longtable}{lp{4.5cm}p{6cm}}
\caption{Módulos excluidos y justificación}
\label{tab:modulos-excluidos}\\
\toprule
\textbf{Módulo / Actividad} & \textbf{Motivo} & \textbf{Justificación técnica}\\
\midrule
\endfirsthead
\toprule
\textbf{Módulo / Actividad} & \textbf{Motivo} & \textbf{Justificación técnica}\\
\midrule
\endhead
\bottomrule
\endfoot

Apps Electron / Tauri & Fuera del sprint & Las apps de escritorio comparten el mismo código web; los defectos de la app de escritorio son reflejo de los del frontend web. Probar ambas en el plazo de 4 días sería redundante.\\

Pruebas de carga (>2h) & Restricción de tiempo & Un ciclo de prueba de rendimiento bajo carga requiere un entorno dedicado y al menos 8h de ejecución. Se documenta como riesgo (véase Cap.~\ref{ch:riesgos}).\\

Automatización completa de suite & Restricción de recursos & El equipo SQA tiene 4 personas y 4 días. Se priorizan scripts de prueba sobre los 2 casos de mayor complejidad (CP-07, CP-08). La automatización completa queda como recomendación.\\

Corrección de código & Fuera del rol SQA & El equipo SQA reporta defectos; la corrección es responsabilidad del equipo de desarrollo. El SQA solo verifica la corrección mediante re-test.\\

Módulo de Administración (Admin) & Bajo riesgo & El panel de administración es de uso interno y su impacto en el usuario final es indirecto. No presenta flujos críticos de negocio.\\

Módulo BNR (ruido de fondo) & Dependencia de hardware & La eliminación de ruido de fondo (Background Noise Removal vía gRPC) requiere hardware específico de NVIDIA Riva. No disponible en el entorno de pruebas local.\\

\bottomrule
\end{longtable}

% ---------------------------------------------------------------------------
\section{Resumen del Sistema}
\label{sec:resumen-sistema}

SpeechNotes implementa un pipeline de transcripción de audio a texto en cinco etapas:

\begin{enumerate}
    \item \textbf{Captura de audio:} El navegador captura el micrófono mediante WebRTC. El VAD (Voice Activity Detection) detecta inicio/fin de habla en tiempo real con Socket.IO.
    \item \textbf{Reducción de ruido (BNR):} Opcional, vía NVIDIA Riva gRPC. Elimina ruido de fondo antes de la transcripción.
    \item \textbf{Transcripción ASR:} El audio limpio se envía a NVIDIA NIM (modelo Parakeet TDT 0.6B para inglés, Canary 1B Flash para español) y se obtiene texto.
    \item \textbf{Formateo con LLM:} Un agente basado en Llama 3.1 70B reestructura el texto crudo en un documento organizado con encabezados y párrafos.
    \item \textbf{Chat / RAG:} El usuario puede interrogar sus transcripciones mediante un agente RAG que usa ChromaDB como almacén vectorial y NVIDIA NIM como LLM.
\end{enumerate}

% ---------------------------------------------------------------------------
\section{Relación con Otros Documentos}
\label{sec:relacion-docs}

Este SQAP se complementa con:
\begin{itemize}
    \item \textbf{Plan de Gestión de Configuración (SCMP):} Define el control de ramas, versionado semántico y etiquetas de entrega.
    \item \textbf{Guía del Proyecto Final SQA:} Parámetros académicos, criterios de evaluación y plazos de entrega.
    \item \textbf{Registro de Defectos Jira:} Base oficial de defectos (BUG-001 a BUG-009) referenciados en los apéndices.
\end{itemize}
